\section{Mejoras al proyecto actual}\label{sec:mejorar_proyecto_actual}

Las siguientes mejoras corresponden a funcionalidades identificadas durante el 
desarrollo que quedaron fuera del alcance de la titulación. Se recomienda abordarlas 
como primera etapa de cualquier proyecto de continuación, ya que complementan 
directamente el módulo de préstamos ya implementado.

\subsection{Seguridad y autenticación}

\subsubsection*{Refresh Tokens}

La implementación actual obliga al usuario a volver a iniciar sesión cada vez que el 
token JWT expira. Se propone implementar un mecanismo de refresh tokens que renueve 
automáticamente la sesión sin interrumpir el flujo de trabajo. Esto requiere 
configurar Litestar para emitir refresh tokens junto al access token, y modificar el 
interceptor de \texttt{apiFetch} en el frontend para intentar el refresco automático 
ante un error \texttt{401} antes de redirigir al login.

\subsubsection*{Expiración automática de cuentas}

El campo \texttt{is\_active} actualmente es controlado de forma manual por el 
administrador. Una mejora sería vincularlo a reglas automáticas, como desactivar 
cuentas que no han tenido actividad durante un período definido, o asignar fecha de 
vencimiento a cuentas temporales de alumnos de semestres específicos.

\subsection{Módulo de Préstamos}

\subsubsection*{Cancelación de solicitudes por el usuario}

Actualmente el usuario no puede cancelar una solicitud una vez enviada. Se propone 
permitir la cancelación mientras la solicitud esté en estado \texttt{"pendiente"}, 
antes de que el administrador la revise.

\subsubsection*{Catálogo agrupado por producto}

La vista actual del catálogo muestra una fila por cada dispositivo físico individual, 
lo que expone códigos internos que no son relevantes para el usuario. Se propone 
agrupar los dispositivos disponibles por producto, mostrando solo el nombre del 
producto con un contador de unidades disponibles. La asignación del dispositivo 
físico específico quedaría a criterio del administrador al registrar la entrega.

\subsubsection*{Serialización correcta del DTO de préstamos}

La relación \texttt{device.product} dentro de los \texttt{loan\_request\_items} no 
se serializa correctamente desde el DTO del backend. Actualmente se resuelve con 
llamadas adicionales desde el frontend. Se propone resolver este problema a nivel del 
\texttt{LoanRequestReadDTO} para eliminar las llamadas redundantes.

\subsection{Módulo de Dispositivos}

\subsubsection*{\texttt{DeviceForm} reutilizable}

El formulario de creación de dispositivos está embebido en \texttt{DevicesListView}. 
Se propone extraerlo a un componente independiente \texttt{DeviceForm.vue} para 
poder reutilizarlo desde \texttt{ProductDetailView} mediante un \texttt{Dialog} con 
el \texttt{productId} precargado.

\subsubsection*{Botones de \texttt{ProductDetailView}}

El botón ``Editar producto'' actualmente navega a la lista de productos en lugar de 
abrir el Drawer de edición directamente. El botón ``Agregar dispositivo'' muestra un 
mensaje de próximamente en lugar de abrir un formulario. Ambos deben ser 
implementados correctamente.

\subsection{Módulo de Usuarios}

\subsubsection*{Drawer para ver detalles de usuario}

El botón ``Ver detalles'' en \texttt{UsersListView} abre un Dialog simple. Se propone 
reemplazarlo por un Drawer lateral con el mismo diseño que el Drawer de detalle de 
préstamos, mostrando toda la información del usuario de forma más cómoda y 
consistente.

\subsubsection*{Perfil personalizable en \texttt{MyUserView}}

Se propone extender la vista de perfil personal con dos funcionalidades: una foto de 
perfil de tamaño reducido (almacenada como \texttt{base64} o referencia de archivo 
con límite de tamaño), y un sistema de secciones académicas personalizables donde el 
usuario puede completar campos como descripción personal, proyectos desarrollados, 
asignaturas cursando, aprobadas, favoritas y no favoritas con justificación. Estas 
secciones podrían almacenarse como un campo JSON en el modelo \texttt{User} del 
backend para mayor flexibilidad.

% -------------------------------------------------------
\section{Nuevos módulos propuestos}

Los siguientes módulos representan extensiones naturales de la plataforma hacia otras 
necesidades del LabDIC. Cada uno puede desarrollarse de forma independiente como 
proyecto de titulación, conectándose a la infraestructura existente.

\subsection{Módulo 1 --- Préstamo de Inventario (continuación)}

Corresponde a la consolidación de todas las mejoras pendientes descritas en la 
sección \ref{sec:mejorar_proyecto_actual}. Incluye la resolución de los problemas técnicos identificados, la 
implementación de las funcionalidades que quedaron fuera del alcance y la 
incorporación de mejoras de UX identificadas durante el desarrollo.

\subsubsection*{Alcance sugerido}

\begin{itemize}
	\item Todas las mejoras de la sección anterior de este documento.
	\item Implementación de refresh tokens y expiración automática de cuentas.
	\item Catálogo de dispositivos agrupado por producto.
	\item Perfil de usuario con foto y secciones académicas.
	\item Cancelación de solicitudes de préstamo por el usuario.
	\item Optimización de DTOs del backend para eliminar llamadas redundantes.
\end{itemize}

\subsection{Módulo 2 --- Administración de Inventario}

Módulo orientado a tareas administrativas del laboratorio que van más allá de los 
préstamos. Toma como base el modelo de datos ya desarrollado (productos, dispositivos, 
ubicaciones) y agrega herramientas de gestión documental y reportería para los 
administradores del LabDIC.

\subsubsection*{Alcance sugerido}

\begin{itemize}
	\item Dashboard administrativo con estadísticas del inventario: dispositivos por 
	estado, por ubicación y por categoría.
	\item Exportación de reportes en formatos \texttt{.xlsx}, \texttt{.pdf} y 
	\texttt{.docx} con el inventario completo o filtrado.
	\item Generación de documentos internos administrativos: memorándum de baja de 
	dispositivos y memorándum de traslado de dispositivos entre ubicaciones.
	\item Historial completo de movimientos de cada dispositivo con trazabilidad por 
	usuario.
	\item Alertas de dispositivos en mantención prolongada o con préstamos vencidos.
	\item Gestión de proveedores y registro de adquisiciones de equipos.
\end{itemize}

\subsection{Módulo 3 --- Gestión de DHCP}

Los administradores del LabDIC gestionan manualmente el archivo de configuración DHCP 
para controlar qué dispositivos externos pueden conectarse a la red WiFi del 
laboratorio. Este módulo propone una interfaz que abstraiga la manipulación directa 
del archivo DHCP, mejorando significativamente la experiencia del administrador y 
reduciendo errores de configuración.

\subsubsection*{Contexto técnico}

El sistema DHCP manual del laboratorio asigna direcciones IP fijas a dispositivos 
identificados por su dirección MAC. Las direcciones están organizadas en rangos IPv4 
según el tipo de dispositivo y la sala a la que pertenecen: dispositivos inalámbricos 
temporales de estudiantes, equipos fijos de salas TI-1, TI-2, Redes, entre otros.

\subsubsection*{Alcance sugerido}

\begin{itemize}
	\item Interfaz web para agregar, editar y eliminar registros DHCP sin editar el 
	archivo de configuración directamente.
	\item Formulario de registro con: nombre de usuario, nombre del dispositivo, 
	dirección MAC, tipo de dispositivo y sala o rango IPv4 de destino.
	\item Validación automática de MAC duplicadas antes de agregar un nuevo registro.
	\item Control de cuota por tipo de usuario: máximo de dispositivos configurables 
	por rol (por ejemplo, 3 para usuario común, ilimitados para profesores y 
	administradores).
	\item Asignación de fecha de expiración para dispositivos temporales de 
	estudiantes por semestre, con eliminación automática al vencer.
	\item Visualización del estado actual del archivo DHCP organizado por rangos IPv4 
	y salas.
	\item Exportación del archivo de configuración DHCP listo para aplicar en el 
	servidor.
\end{itemize}

Este módulo tiene complejidad técnica significativa por su interacción con la 
infraestructura de red del laboratorio. Se recomienda como proyecto de titulación 
independiente con participación del equipo de administración de redes del 
departamento.

\subsection{Módulo 4 --- Administración de Salas}

Módulo de gestión de horarios y reservas de las salas de clases del LabDIC. Permite 
a los profesores solicitar salas, visualizar disponibilidad en formato calendario y 
gestionar cambios de sala, todo bajo supervisión y aprobación de los administradores 
del laboratorio.

\subsubsection*{Contexto}

El LabDIC administra varias salas de clases: TI-1, TI-2, Redes, Magister, entre 
otras. La coordinación de horarios y asignación de salas actualmente se realiza de 
forma manual, lo que puede generar conflictos y falta de visibilidad para los 
profesores.

\subsubsection*{Alcance sugerido}

\begin{itemize}
	\item Vista de calendario por sala con columnas por día de la semana (lunes a 
	viernes) y filas por bloque horario de 1 hora 30 minutos desde las 8:00 hasta 
	las 19:20.
	\item Un calendario independiente por cada sala: TI-1, TI-2, Redes, Magister y 
	otras que se agreguen.
	\item Nuevo rol ``Profesor'' con permisos específicos: visualizar calendarios, 
	solicitar asignación de sala y solicitar cambios.
	\item Flujo de asignación de sala: el profesor solicita una sala para un horario 
	específico, el administrador aprueba o rechaza con motivo.
	\item Solicitud de cambio de sala entre profesores (\textit{swap}): ambos 
	profesores deben aceptar el intercambio antes de que el administrador lo apruebe.
	\item Solicitud de sala temporal por período menor a un mes para actividades 
	extraordinarias.
	\item Vista de administrador con control total: puede asignar, modificar y 
	eliminar cualquier reserva.
	\item Notificaciones o alertas cuando se aprueba o rechaza una solicitud.
\end{itemize}

Este módulo tiene alta complejidad de UX por la naturaleza del calendario interactivo 
y los múltiples flujos de aprobación. Se recomienda como proyecto de titulación 
independiente con participación directa de los profesores del departamento como 
usuarios de prueba.

% -------------------------------------------------------
\section{Resumen y recomendaciones}

La tabla~\ref{tab:trabajo_futuro_resumen} resume los cuatro módulos propuestos con su 
clasificación de complejidad y el tipo de proyecto que mejor se adapta a cada uno.

\begin{table}[H]
	\centering
	\begin{tabular}{p{0.30\linewidth} p{0.22\linewidth} p{0.38\linewidth}}
		\hline
		\textbf{Módulo} & \textbf{Tipo} & \textbf{Recomendación} \\
		\hline
		Préstamo de Inventario (continuación) 
		& Continuación directa 
		& Primera etapa para cualquier proyecto de continuación. \\
		\hline
		Administración de Inventario 
		& Extensión nueva 
		& Puede combinarse con el módulo 1 en una sola titulación. \\
		\hline
		Gestión de DHCP 
		& Módulo nuevo 
		& Titulación independiente con foco en infraestructura de red. \\
		\hline
		Administración de Salas 
		& Módulo nuevo 
		& Titulación independiente con foco en UX y flujos de aprobación. \\
		\hline
	\end{tabular}
	\caption{Resumen de módulos de trabajo futuro propuestos.}
	\label{tab:trabajo_futuro_resumen}
\end{table}

Se recomienda que el primer proyecto de continuación aborde en conjunto los módulos 1 
y 2, consolidando el trabajo existente y añadiendo las capacidades administrativas de 
reportería. Los módulos 3 y 4, por su complejidad técnica y su alcance, se recomiendan 
como proyectos de titulación independientes, idealmente con participación activa del 
personal del LabDIC desde las etapas tempranas de análisis de requerimientos.

La arquitectura modular de la plataforma (Vue~3 + Litestar + PostgreSQL) fue 
diseñada deliberadamente para facilitar la incorporación de nuevos módulos sin 
reestructurar lo existente. Cada módulo nuevo se integra como un conjunto de rutas, 
controladores, servicios y vistas independientes que conviven con los módulos 
actuales.